\section{Informational Efficiency Exercise}

The analytical regions are useful because they tell an experimenter what can be inferred after observing a report.
This section reports a design exercise generated by the reproducible script \texttt{scripts/design\_efficiency.py}.
The exercise draws latent beliefs
\[
p\sim\operatorname{Dirichlet}(\alpha,\ldots,\alpha)
\]
and, for each rule, computes an optimal report \(r_S(p)\) and the bounds induced by \(P_S(r_S(p))\).
The final run uses \(5{,}000\) belief draws for each cell in
\[
n\in\{5,10,20,50\},\qquad
k\in\{2,3,5,10\},\qquad
\alpha\in\{0.1,0.3,1,3,10\}.
\]
Manhattan bounds are threshold-computed with tolerance \(10^{-4}\).

The comparison is conditional on the theoretical characterizations above.
It is not evidence of incentive compatibility, and it is not an optimality theorem.
It is a diagnostic for the practical question: which rule gives a researcher tighter inference for the object of interest?

\paragraph{Aggregate comparison.}
Tables \ref{tab:design-widths} and \ref{tab:design-wins} report averages over the full design grid.
Lower widths indicate more precise inference.
Win shares report the average share of belief draws for which a rule gives the narrowest interval for the stated target.

\begin{table}[h]
\centering
\small
\begin{tabular}{lrrrr}
\hline
Rule & avg coord & max coord & linear mean & skewed mean\\
\hline
Discrete metric & 0.1025 & 0.1478 & 0.1142 & 0.1025\\
Quadratic distance & 0.1042 & 0.1185 & 0.1112 & 0.1042\\
Manhattan distance & 0.1021 & 0.1253 & 0.1101 & 0.1021\\
\hline
\end{tabular}
\caption{Average widths in the latent-belief design exercise. Columns average over all \((n,k,\alpha)\) cells. Manhattan entries are threshold-computed.}
\label{tab:design-widths}
\end{table}

\begin{table}[h]
\centering
\small
\begin{tabular}{lrrr}
\hline
Rule & avg-coordinate wins & max-coordinate wins & linear-mean wins\\
\hline
Discrete metric & 0.6204 & 0.3431 & 0.5266\\
Quadratic distance & 0.3189 & 0.6070 & 0.3776\\
Manhattan distance & 0.0607 & 0.0499 & 0.0958\\
\hline
\end{tabular}
\caption{Average win shares in the latent-belief design exercise. A win is the smallest interval width among the three rules for a draw and target. The linear mean uses \(x_i=(i-1)/(k-1)\).}
\label{tab:design-wins}
\end{table}

The aggregate results do not support a universal ranking.
Discrete-metric scoring gives the narrowest average coordinate interval most often.
Quadratic-distance scoring is strongest for worst-coordinate uncertainty, with a win share of about \(0.607\) for maximum coordinate width.
It is also competitive for mean inference and performs especially well in sparse-belief environments.
Manhattan distance is a useful structured comparison rule, but this run does not support presenting it as the dominant design choice.

\paragraph{Where quadratic distance helps.}
The cell-level results show that quadratic distance is particularly useful when beliefs are sparse or concentrated.
For example, with \(\alpha=0.1\), quadratic distance wins the linear-mean criterion in more than \(0.8\) of draws for many cells, including \(n=50,k=10\), \(n=50,k=5\), and \(n=20,k=5\).
It also wins the worst-coordinate criterion almost always in several larger-\(k\) cells.
This is the setting in which the projection interpretation is most valuable: the rule translates reports near sparse or extreme count vectors into tight bounds on high-probability categories and mean-oriented functionals.

\paragraph{Payment probability.}
The discrete-metric rule has a fixed-prize implementation, which is attractive under risk aversion.
Its implementation cost is that the exact-match payment probability can be small.
In the final run, the average exact-match probability is high for very sparse binary beliefs, but it falls sharply as \(n\), \(k\), and belief balance increase; for instance, it is essentially zero in the \(n=50,k=10,\alpha\in\{1,3,10\}\) cells.
Thus narrow bounds under discrete-metric scoring should be weighed against the chance that the subject rarely wins the fixed prize.

\paragraph{Fixed-report illustrations.}
Table \ref{tab:selected-design} holds the observed report fixed and compares the induced bounds.
This answers a different question from the latent-belief simulation: conditional on seeing the same report, how does the scoring rule change the inferred belief set?

\begin{table}[h]
\centering
\small
\begin{tabular}{llrrr}
\hline
Report and rule & payment prob. & avg coord width & linear mean interval & method\\
\hline
\multicolumn{5}{l}{\(n=10,\ k=3,\ r=(5,3,2)\)}\\
Discrete metric & 0.0850 & 0.1162 & \([0.3182,0.4167]\) & closed form / LP\\
Quadratic distance &  & 0.1333 & \([0.3000,0.4000]\) & closed form / LP\\
Manhattan distance &  & 0.1278 & \([0.3055,0.4050]\) & threshold\\
\hline
\multicolumn{5}{l}{\(n=2,\ k=3,\ r=(1,1,0)\)}\\
Discrete metric & 0.5000 & 0.3889 & \([0.1667,0.5000]\) & closed form / LP\\
Quadratic distance &  & 0.5000 & \([0.1250,0.5000]\) & closed form / LP\\
Manhattan distance &  & 0.4492 & \([0.1464,0.5000]\) & threshold\\
\hline
\multicolumn{5}{l}{\(n=20,\ k=5,\ r=(10,4,3,2,1)\)}\\
Discrete metric & 0.0061 & 0.0714 & \([0.2381,0.3152]\) & closed form / LP\\
Quadratic distance &  & 0.0800 & \([0.2125,0.2875]\) & closed form / LP\\
Manhattan distance &  & 0.0770 & \([0.2182,0.2948]\) & threshold\\
\hline
\end{tabular}
\caption{Fixed-report illustrations. The payment probability is the discrete-metric exact-match probability at \(p=r/n\); it is an implementation cost of that fixed-prize rule, not a property of the other distance rules.}
\label{tab:selected-design}
\end{table}

The fixed-report cases again show why the paper should not claim uniform dominance.
For some reports, discrete-metric scoring gives narrower average coordinate intervals.
For others, quadratic distance gives narrower mean bounds.
The relevant design choice depends on whether the experimenter prioritizes average coordinate precision, worst-coordinate precision, mean inference, implementation robustness, or ease of explanation.
